\section{Introduction}
\subsection{Scientific/Technological Context}

In the rapidly evolving landscape of technology and artificial intelligence, Computer Vision has transitioned from a niche research area to a cornerstone of modern applications, namely Human-Computer Interaction (HCI)~\cite{kumari2025hand}.
At its core lies Image Processing, a fundamental technique that enables computers to interpret and manipulate meaningful data from visual inputs.

The technology to bridge the gap between human intent and machine execution through real-time hand gesture recognition is increasingly relevant in various domains, such as sterile medical environments, immersive augmented reality and assistive tools for individuals with motor impairments~\cite{mopidevi2023hand}.
By moving beyond traditional input devices we can create more intuitive and natural interfaces that enhance user experience and accessibility.

Technologically this lab script centers on the synergy between OpenCV, a powerful open-source computer vision library and Google's MediaPipe, a framework for building multimodal applied machine learning pipelines, specifically its hand tracking solution that utilizes a pre-trained palm detector and a hand landmark model to provide a 21-point 3D hand skeletal map~\cite{biswas2024mediapipe, mediapipe_hands}.

\subsection{Motivation for this topic/lab script}

In large part due to the COVID-19 pandemic, the world has witnessed a significant shift towards remote interactions and contactless technologies.
This has accelerated the demand for more intuitive and natural interfaces that can bridge the gap between humans and machines.

Coupled with the advancements in machine learning and the availability of powerful libraries such as OpenCV and MediaPipe, tackling the challenge of real-time hand gesture recognition seemed like a natural choice for a lab script in the context of Advanced Computing Systems (SISTCA).

\subsection{Objectives (of this lab script)}

The primary objective of this lab script is the development of a distributed real-time hand gesture recognition system. This document defines both the theoretical concepts to be acquired and the practical milestones to be achieved.

Upon completion, students should be able to:
\begin{itemize}
    \item \textbf{Learn (Theoretical Concepts):} Understand the principles of edge computing, spatial anatomical feature extraction (landmarks), and the fundamentals of lightweight Local Area Network (LAN) communication through socket protocols.
    \item \textbf{Do (Practice and Milestones):}
    \begin{itemize}
        \item Configure a Python development environment.
        \item Capture video streams and extract hand topology using the MediaPipe library.
        \item Implement geometric logic to identify and count raised fingers.
        \item Develop socket communication to transmit the gesture classification to another machine on the network with minimum latency.
    \end{itemize}
\end{itemize}

\subsection{Structure (of this lab script)}

The remainder of this document is organized as follows:
\begin{itemize}
    \item Section 2 (Theoretical): Presents the fundamental theoretical concepts and the state-of-the-art of the addressed technologies.
    \item Section 3 (Outlook of the technology): Describes the hardware and software configuration requirements to start the project.
    \item Section 4 (Tutorial/Functionality): Constitutes the core of this lab script, guiding the student through a step-by-step tutorial to build the system.
    \item Section 5 (Exercises): Proposes practical exercises and provides their respective resolutions.
    \item Section 6 (Challenge): Presents an autonomous challenge for the students to consolidate the learned concepts.
    \item Section 7 (References): Lists all the bibliographic references used throughout the document.
    \item \textbf{Annexes:} Contains troubleshooting resources, including a direct link to a repository where you can download a dataset and a pre-trained model.
\end{itemize}